% Generated draft: requires mathematical and source review.
\documentclass[11pt]{article}
\usepackage[margin=25mm]{geometry}
\usepackage{fontspec}
\setmainfont{Latin Modern Roman}
\usepackage{amsmath,amssymb,mathtools,unicode-math}
\setmathfont{Latin Modern Math}
\usepackage{enumitem,needspace,xurl,hyperref}
\setlength{\parindent}{0pt}
\setlength{\parskip}{5pt}
\setlength{\emergencystretch}{3em}
\newcommand{\N}{\mathbb{N}}
\newcommand{\Z}{\mathbb{Z}}
\newcommand{\Q}{\mathbb{Q}}
\newcommand{\R}{\mathbb{R}}
\newcommand{\C}{\mathbb{C}}
\newcommand{\F}{\mathbb{F}}
\newcommand{\A}{\mathbb{A}}
\newcommand{\PP}{\mathbb{P}}
\newcommand{\E}{\mathbb{E}}
\newcommand{\eps}{\varepsilon}
\newcommand{\dd}{\,\mathrm{d}}
\begin{document}
\section*{UnsolvedMath ID: 701307}
\textbf{Source problem code:} catalogue entry
\subsection{Definitions and mathematical statement}
For a group \(G\), use the commutator convention \([x,y]=x^{-1}y^{-1}xy\) and left-normed iterated commutators \([x,y,z]=[[x,y],z]\). An element \(a\in G\) is called a left \(3\)-Engel element if
\[
[x,a,a,a]=1\qquad\text{for every }x\in G.
\]
Define the lower central series by \(\gamma_1(G)=G\) and
\[
\gamma_{n+1}(G)=\langle [u,v]:u\in\gamma_n(G),\ v\in G\rangle.
\]
The group \(G\) is nilpotent of class at most \(c\) if \(\gamma_{c+1}(G)=\{1\}\).

The question is whether there exists a function
\[
f:\Z_{\geq 1}\longrightarrow\Z_{\geq 0}
\]
such that, for every positive integer \(d\), every nilpotent group \(G\), and every generating tuple \(a_1,\ldots,a_d\) satisfying
\[
G=\langle a_1,\ldots,a_d\rangle,
\qquad
[x,a_i,a_i,a_i]=1
\quad(x\in G,\ 1\leq i\leq d),
\]
one has \(\gamma_{f(d)+1}(G)=\{1\}\). No finiteness hypothesis on \(G\) is stated in the source, and the generating tuple is not required to be minimal.

As background, the supplied record states that every group generated by two left \(3\)-Engel elements is nilpotent of class at most \(4\), citing \emph{J. Pure Appl. Algebra} \textbf{188} (2004), 1--6. This background assertion and its incomplete bibliographic reference are reproduced as supplied; they have not been independently verified.
\subsection{Short English statement}
Does the nilpotency class of a nilpotent group generated by \(d\) left \(3\)-Engel elements admit a bound depending only on \(d\)? The supplied record attributes the problem to A. Abdollahi and reports a class bound of \(4\) for groups generated by two such elements.
\subsection{Sources}
\begin{itemize}
\item Local UnsolvedMath registry, permanent ID \(701307\), exact source code \texttt{catalogue entry}, title: \emph{Kourovka Notebook Problem 17.5}. The registry labels the problem open; this is supplied source metadata, not an independently verified assessment.
\item \emph{The Kourovka Notebook (Unsolved Problems in Group Theory)}, original problem 17.5, proposed by A. Abdollahi in the 17th issue (2010). The supplied background also identifies the 21st issue edition (2026), arXiv:1401.0300v48. The collection description names E. I. Khukhro and V. D. Mazurov as editors.
\item Supplied reference for the two-generator background assertion: \emph{J. Pure Appl. Algebra} \textbf{188} (2004), 1--6. The record supplies neither the article title nor its authors.
\end{itemize}
\end{document}
